\section{Early collisions in the circular unitary ensemble}
\label{sec:cue-collisions}

We apply the local dynamics to the extreme gaps whose distribution was proved by Ben Arous and Bourgade. Let $U_N$ be Haar distributed in $U(N)$ and write its ordered eigenangles as
$0\le\theta_1<\cdots<\theta_N<2\pi$, with
$\theta_{i+N}=\theta_i+2\pi$. Define
\[
 Z_N(z)=\prod_{i=1}^N2\sin\frac{z-\theta_i}{2},\qquad
 Z_{N,-a}=e^{aD^2}Z_N,\qquad
 \delta_N=\min_i(\theta_{i+1}-\theta_i).
\]
The divisor of $Z_N$ is periodic, including when $Z_N$ itself is antiperiodic. Its first loss of reality occurs at
$A_N=\inf\{a>0:Z_{N,-a}\text{ has a nonreal zero}\}$.

\begin{theorem}[First collision and the early collision process]\label{thm:cue-collisions}
For $N\ge2$, almost surely $0<A_N<\infty$. Moreover,
\[
 \frac{8A_N}{\delta_N^2}=1+O_{\mathbb P}(N^{-1/3}),\qquad
 \lim_{N\to\infty}\mathbb P\left(
   \frac{8N^{8/3}A_N}{(72\pi)^{2/3}}>v\right)=e^{-v^{3/2}}
 \quad(v\ge0).
\]
With probability tending to one the first colliding pair starts at the smallest cyclic gap.
For each real double collision, let $a_{N,j}>0$ be its negative-time magnitude and $b_{N,j}$ its circular location. On bounded rescaled time intervals the point process
\[
 \sum_j\delta_{(8N^{8/3}a_{N,j}/(72\pi)^{2/3},\,b_{N,j})}
\]
converges to the Poisson process with intensity
\[
                  \frac32\sqrt v\,dv\,\frac{db}{2\pi}.
\]
With probability tending to one, every collision on such an interval is double, creates one nonreal conjugate pair, and that pair remains nonreal for the rest of the interval. In particular the number of nonreal pairs at rescaled time $v$ converges to $\operatorname{Poisson}(v^{3/2})$.
\end{theorem}

\begin{proof}
We first control the entire finite divisor. Write
$Z_N(z)=\sum_{k=0}^Nc_ke^{i(k-N/2)z}$. Its endpoint coefficients $c_0,c_N$ are nonzero. Heat multiplication changes $c_k$ to $c_ke^{-a(k-N/2)^2}$. Multiplication by $e^{iNz/2}$ gives a polynomial of degree $N$ in $w=e^{iz}$ with nonzero constant coefficient. On each compact interval of $a$, its roots lie in a compact annulus: the usual polynomial root bound applies to the polynomial and its reciprocal. A zero therefore cannot leave the real axis by escape to imaginary infinity at a finite time. A simple real zero stays real by conjugation and the implicit function theorem. The first loss of reality must occur at a real collision.

Before that time, the real angles $y_i(a)$ satisfy
\[
 y_i'(a)=-\sum_{j\ne i}\cot\frac{y_i(a)-y_j(a)}2.
\]
This follows by differentiating $Z_{N,-a}(y_i(a))=0$ and its sine product. On an empty adjacent arc of length $g=y_{i+1}-y_i$, the mutual terms give $-2\cot(g/2)$ to $g'$. Each other cotangent difference is negative before the overall minus sign, since the arc contains no pole. Thus
\[
          g'\ge-2\cot(g/2),\qquad (g^2)'\ge-8.
\]
Equivalently,
\[
                  (\cos(g/2))'\le\cos(g/2).
\]
Hence a closing gap has closing time at least $-\log\cos(g(0)/2)$.
A nonpositive initial cosine rules out closing by this inequality.
Since $\delta_N<\pi$ almost surely,
\[
                   A_N\ge-\log\cos(\delta_N/2)\ge\delta_N^2/8.
\]
Positivity follows. For even $N$, the constant Fourier coefficient is almost surely nonzero and eventually dominates, leaving no real zeros. For odd $N$, the nonzero frequencies $\pm1/2$ eventually dominate and leave one simple real zero per period. The relevant coefficients are nonzero real analytic functions of the angles: this can be checked when all angles tend to zero, where the product tends to $(2\sin(z/2))^N$. Their zero sets have Lebesgue measure zero, hence zero CUE probability. Uniform $C^1$ convergence after the appropriate rescaling proves the asserted eventual real-zero counts. This proves finiteness of $A_N$.

Put $\ell_N=N^{-7/6}$. With probability tending to one, no circular arc of length $\ell_N$ contains three eigenangles. Indeed, with
$v(x)=(1,e^{ix},\ldots,e^{i(N-1)x})/\sqrt{2\pi}$, the three-point intensity is
$\|v(x_1)\wedge v(x_2)\wedge v(x_3)\|^2$.
Successive divided differences extract $\prod_{i<j}(x_i-x_j)$, and the norms of $v,v',v''$ are respectively $O(N^{1/2}),O(N^{3/2}),O(N^{5/2})$. Hence
\[
 \rho_3(x_1,x_2,x_3)\le CN^9\prod_{i<j}|x_i-x_j|^2.
\]
Integrating over triples in a forward arc of length $\ell_N$ gives expected count at most $CN^9\ell_N^8=O(N^{-1/3})$. Every offending arc has such a triple with its first point as the initial endpoint. Markov's inequality proves the claim.

Ben Arous--Bourgade \cite[Theorem~1.1, p.~2650; Corollary~1.2, p.~2651]{ben-arous-bourgade} prove that the marked gaps $(N^{4/3}\delta,\theta)$ converge to a Poisson process with intensity
$u^2du/(24\pi)\,d\theta/(2\pi)$. In particular
\[
 \frac{N^{4/3}\delta_N}{(72\pi)^{1/3}}\Longrightarrow Y,
 \qquad \mathbb P(Y>y)=e^{-y^3},
\]
and the ratio of the second-smallest to smallest gap converges to a variable strictly larger than one almost surely.
On the no-triples event, once $\delta_N\le\ell_N/2$, every complementary periodic zero of the minimum pair is at distance at least $\ell_N/2$ from its midpoint. At most two zeros lie in each interval of length $\ell_N$, so summing the inverse-square shells gives $S_N\le C\ell_N^{-2}$. Consequently
$\kappa_N=\delta_N^2S_N=O_{\mathbb P}(N^{-1/3})$.
Theorem~\ref{thm:isolated-pair} gives a collision at
$a_{\min}=\delta_N^2(1+\kappa_N/4+O(\kappa_N^2))/8$.
The global lower bound and $A_N\le a_{\min}$ prove the ratio estimate. With probability tending to one, the square of every other initial gap divided by $8$ exceeds $a_{\min}$, by the joint extreme-gap limit. The differential inequality therefore identifies the first pair. Squaring the limit variable $Y$ proves the stated Weibull law.

To pass all early collisions, including those after the first one, fix $A>0$ and set
\[
 s_N=N^{-4/3},\quad T_N=As_N^2,\quad
 B_N=N^{1/24},\quad h_N=B_Ns_N.
\]
On the no-triples event, call every gap $\delta\le h_N$ a pair and every remaining point a singleton. For large $N$, $2h_N<\ell_N$, so the pairs are disjoint. A pair is at distance at least $\ell_N-h_N$ from every other zero. Singletons are mutually at distance at least $h_N$.

We record explicitly the uniform local estimate used for this partition. After removing a singleton at $b$, scale the exact conjugated Gaussian transform by $\sqrt T$, write $\tau=a/T$, and obtain
\[
 I(\tau,u)=\mathbb E[Y R(\sqrt T Y)]=u+O(TS),
 \quad Y=u+\sqrt{2\tau}V.
\]
For a pair use $R_0^2=\delta^2+T$, $d=\delta/R_0\le1$, and $\tau=a/R_0^2\le1$ to obtain
\[
 J(\tau,u)=\mathbb E[(Y^2-d^2/4)R(R_0Y)]
           =u^2+2\tau-d^2/4+O(R_0^2S).
\]
The bounds of Lemma~\ref{lem:centered-product} dominate every remainder and every fixed derivative on bounded complex sets, uniformly in $0<d\le1$. For sufficiently small $TS$ or $R_0^2S$, Rouch\'e on $|u|=2$ gives exactly one or two roots, respectively. The singleton is simple and real. In the pair case $J_u$ has one zero, $J_{uu}=2+O(R_0^2S)>0$ on the real disk, and the critical value has derivative $2+O(R_0^2S)>0$. Its initial value is negative because the initial two roots are $\pm d/2$. Thus the pair has at most one collision and cannot return to the real axis during the interval. When its collision lies in the interval, Theorem~\ref{thm:isolated-pair} gives its time, by uniqueness of continuation from the same initial two roots. This establishes the uniform estimate even when $\delta^2/T$ tends to zero.

For each pair, shell summation and the cotangent product give
\[
 S_{\rm pair}\le C\ell_N^{-2},\qquad
 |L_{\rm pair}|\le C\ell_N^{-1}\log(2/\ell_N).
\]
For each singleton they give
\[
 S_{\rm single}\le C(h_N^{-2}+\ell_N^{-2}),\qquad
 |L_{\rm single}|\le C(h_N^{-1}+\ell_N^{-1}\log(2/\ell_N)).
\]
The inverse-square estimates include all periodic copies. For the first logarithmic derivative use the finite cotangent sum, bounded by the harmonic sum of the circular distance shells; the two pair cotangents cancel at the midpoint. These estimates imply
\[
 (\delta^2+T_N)S_{\rm pair}=O(N^{-1/4}),\qquad
 T_NS_{\rm single}=O_A(B_N^{-2}+N^{-1/3}).
\]
Hence the uniform local estimates apply simultaneously to all clusters.
The moving pair disks have radii $O_A(h_N)=o(\ell_N)$ and center shifts $O_A(T_N\ell_N^{-1}\log(2/\ell_N))=o(\ell_N)$. The singleton disks have radius $O_A(s_N)$ and center shifts
$O_A(T_N/h_N+T_N\ell_N^{-1}\log(2/\ell_N))=o(s_N)$.
They remain disjoint modulo $2\pi$. Their root counts add to $N$, and the polynomial in $w$ has degree $N$ throughout, so these disks exhaust every root at every time in $[0,T_N]$.

For the pairs,
$a_c=\delta^2(1+O(N^{-1/4}))/8$ uniformly. Only gaps $\delta=O_A(s_N)$ can collide in this interval; their sharper interaction bound is $\delta^2S=O_A(N^{-1/3})$. Their collision positions differ from their initial left endpoints by $o(1)$. The early collision process is therefore the image of the extreme-gap process, up to uniformly vanishing displacements, under
\[
 (y,b)\longmapsto(y^2,b),\qquad
 y=N^{4/3}\delta/(72\pi)^{1/3}.
\]
The source intensity becomes $3y^2dy\,db/(2\pi)$ in this coordinate and then $(3/2)\sqrt v\,dv\,db/(2\pi)$ after squaring. Tightness of counts in bounded windows follows from the cited point-process convergence, and the limiting process has no atoms at deterministic boundaries. This justifies the displacement comparison and gives all asserted count limits.
\end{proof}

For zeta zeros, Theorem~\ref{thm:isolated-pair} already applies to any adjacent simple pair with small $\delta^2S$. An arithmetic version of the Poisson conclusion requires a marked extreme-gap limit and quantitative control of the complementary inverse-square sums for those pairs, together with exhaustion of early collisions. The exact restricted Fourier support in the Montgomery and Rudnick--Sarnak theorems controls the pair sums used in Section~\ref{sec:survival}; the empty-interval and extreme-gap quantities in the present proof are additional quantities to estimate. No arithmetic Poisson limit is proved here.
